\subsection{4. Die richtige einfache Kette: Wiederholung von
H}\label{die-richtige-einfache-kette-wiederholung-von-h}

\subsubsection{4.1 Der Fehler war die Gleichsetzung zweier verschiedener
Stufen}\label{der-fehler-war-die-gleichsetzung-zweier-verschiedener-stufen}

Die Vektoren \(v_k=T_+^kF\) sind nach Bosonzahl geordnet. Schon bekannt
ist

\[\adjustbox{max width=\linewidth}{$\displaystyle 
T_-v_3=\frac{299520}{229}v_2+w_2,\qquad
\langle v_2,w_2\rangle=0,\qquad
\|w_2\|^2=\frac{5001523200}{229}>0.
$}\]

Diese Rechnung wurde frisch reproduziert. Aus einer kleinen relativen
Norm dieses Seitenzweigs folgt keine Konvergenz der gesamten
Grundzustands- oder Spektralrechnung. Insbesondere fehlt dann immer noch
der große Übergang nach v4.

Eine Lanczos-Kette entsteht stattdessen aus \(F,HF,H^2F,\ldots\) durch
Orthogonalisierung. Für einen selbstadjungierten endlichen
Hamiltonoperator beschreibt sie dessen zyklischen Teilraum exakt. Das
ist ein zulässiges einfaches Rechenbild, aber keine Behauptung, dass
ihre Indizes physikalische Orte oder Bosonzahlen seien. Auch die gesamte
Bank muss nicht mit diesem einen zyklischen Teilraum übereinstimmen.

\subsubsection{4.2 Die ersten fünf Glieder sind jetzt exakt
bestimmt}\label{die-ersten-fuxfcnf-glieder-sind-jetzt-exakt-bestimmt}

Für die normierte Lanczos-Basis \(u_n\) schreiben wir

\[\adjustbox{max width=\linewidth}{$\displaystyle 
Hu_n=b_nu_{n-1}+a_nu_n+b_{n+1}u_{n+1},\qquad b_n>0.
$}\]

Die ersten Diagonalen lauten

\[\adjustbox{max width=\linewidth}{$\displaystyle 
\frac{(a_0,a_1,a_2,a_3,a_4)}{\Delta}
=\left(0,1,2,3,\frac{105168998}{26292551}\right).
$}\]

Die ersten quadrierten Nebendiagonalen sind

\[\adjustbox{max width=\linewidth}{$\displaystyle 
\frac{(b_1^2,b_2^2,b_3^2,b_4^2)}{g^2}
=\left(480,916,\frac{299520}{229},\frac{78877653}{47632}\right).
$}\]

Bis u3 stimmen die Vektoren, bis auf die Phase bei negativem g, mit den
normierten vk überein. Der nächste Vektor ist proportional zu
\(v_4+w_2\). Dabei liegen v4 und w2 auf verschiedenen Bosonstufen.
Deshalb gilt

\[\adjustbox{max width=\linewidth}{$\displaystyle 
\frac{b_4^2}{g^2}
=\frac{\|v_4\|^2+\|w_2\|^2}{\|v_3\|^2},
\qquad
\langle u_4,N_bu_4\rangle
=\frac{4\|v_4\|^2+2\|w_2\|^2}{\|v_4\|^2+\|w_2\|^2}.
$}\]

Die Bosonzahlvarianz dieses Vektors ist exakt

\[\adjustbox{max width=\linewidth}{$\displaystyle 
\operatorname{Var}_{u_4}(N_b)
=\frac{63416178576}{691298238087601}>0.
$}\]

Die frühere Quotientenfortsetzung mit nur \(\|v_4\|^2/\|v_3\|^2\) ließ
den positiven Zusatz \(1809/47632\) weg. Die korrigierte Kette bleibt
einfach, verliert aber gerade deshalb nicht die wirkliche
Mehrkanalstruktur. Ihre Länge ist nicht durch die 33 möglichen
Bosonzahlen festgelegt.

\subsubsection{4.3 Zehn vollständige Energiemomente der gefüllten
Referenz}\label{zehn-vollstuxe4ndige-energiemomente-der-gefuxfcllten-referenz}

Die Fünfer-Kompression bestimmt \(\mu_n=\langle F,H^nF\rangle\) für n
von 0 bis 9 exakt. Ein Weg, der den nächsten nicht enthaltenen
Lanczos-Vektor erreicht und nach F zurückkehrt, benötigt mindestens zehn
Schritte. Aus der Tridiagonalität folgt deshalb die angegebene
Momentengenauigkeit, ohne die übrige Kette zu erfinden.

Beispielsweise gilt

\[\adjustbox{max width=\linewidth}{$\displaystyle 
\mu_0=1,\quad \mu_1=0,\quad
\mu_2=480g^2,\quad \mu_3=480\Delta g^2,
\quad \mu_4=480\Delta^2g^2+670080g^4.
$}\]

Vollständig lauten die Koeffizienten in
\(\mu_n/\Delta^n=c_2x^2+c_4x^4+c_6x^6+c_8x^8\), mit \(x=g/\Delta\), für
n von 2 bis 9:

{\def\LTcaptype{none} % do not increment counter
\begin{longtable}[]{@{}lrrrr@{}}
\toprule\noalign{}
n & c2 & c4 & c6 & c8 \\
\midrule\noalign{}
\endhead
\bottomrule\noalign{}
\endlastfoot
2 & 480 & 0 & 0 & 0 \\
3 & 480 & 0 & 0 & 0 \\
4 & 480 & 670080 & 0 & 0 \\
5 & 480 & 2219520 & 0 & 0 \\
6 & 480 & 5527680 & 1510510080 & 0 \\
7 & 480 & 12353280 & 10437173760 & 0 \\
8 & 480 & 26213760 & 48253363200 & 4615972423680 \\
9 & 480 & 54144000 & 187581404160 & 54295325184000 \\
\end{longtable}
}

Diese Momente gehören \textbf{F}, nicht dem nativen Grundzustand Omega.
Sie sind auch nicht mit den geladenen Antwortmomenten aus der früheren
Fassung zu verwechseln. Sie machen eine endliche Modellrechnung genauer,
beweisen aber keine physische Präparation von F und keine
Vollständigkeit des Feldwörterbuchs.

\subsubsection{4.4 Kleine, rigorose Verschärfung der
Energiegrenze}\label{kleine-rigorose-verschuxe4rfung-der-energiegrenze}

Am Prüfpunkt g geteilt durch Delta gleich 1/20 hat die niedrigste
Eigenenergie der Fünfer-Kompression die durch rationale Sturm-Zählung
zertifizierte Lage

\[\adjustbox{max width=\linewidth}{$\displaystyle 
-1.12963813<\frac{E_{\mathrm{Ritz},5}}{\Delta}<-1.12963811.
$}\]

Rayleigh-Ritz liefert für den wahren Grundzustand daher

\[\adjustbox{max width=\linewidth}{$\displaystyle 
E_0<-1.12963811\Delta.
$}\]

Die frühere untere Schranke bleibt \(E_0>-1.158089\Delta\). Zusammen mit
dem früheren strengen N=63-Floor \(E_h>-1.121899\Delta\) folgt nun

\[\adjustbox{max width=\linewidth}{$\displaystyle 
\epsilon=E_h-E_0>0.00773911\Delta.
$}\]

Die Verbesserung gegenüber 0.007737 ist klein. Die obere Polgrenze
0.039079764 und das Gewicht über 88.007628 Prozent werden dadurch nicht
ungeprüft verändert. Auch die niedrigste Ritz-Energie ist kein
berechneter Zentralwert von E0.

\subsection{5. Vollständig gelöst: die innere bilineare
Zerlegung}\label{vollstuxe4ndig-geluxf6st-die-innere-bilineare-zerlegung}

Für den 64-dimensionalen Fermionträger gilt

\[\adjustbox{max width=\linewidth}{$\displaystyle 
\operatorname{End}(16\otimes4)
=(1,1)\oplus(45,1)\oplus(210,1)
\oplus(1,15)\oplus(45,15)\oplus(210,15).
$}\]

Die sechs Dimensionen sind 1, 45, 210, 15, 675 und 3150. Der offene Rest
ist also nicht mehr unbestimmt:

\[\adjustbox{max width=\linewidth}{$\displaystyle 
4035=210+675+3150.
$}\]

Eine kurze konkrete Konstruktion genügt. Auf dem chiralen 16er-Raum
liefern Cliffordprodukte der Grade 0, 2 und 4 genau 1, 45 und 210
unabhängige Matrizen. Ihr ganzzahliger Hilbert-Schmidt-Gram ist
\(16I_{256}\). Unter Spin(10)-Konjugation bleiben diese Gradräume
invariant. Auf der Farbseite zerfällt End(4) in die skalare Matrix und
die 15 spurlosen Matrizen. Die Tensorprodukte ergeben die gesamte
Zerlegung. Die Weyl-Multiplizitäten wurden zusätzlich überprüft.

Die verwendete Spinorzerlegung ist etablierte Darstellungstheorie, keine
neue Vorhersage von 210 physikalischen Teilchen. Der Literaturabgleich
bestätigt die drei Spin(10)-Kanäle 1, 45 und 210; unsere Rechnung ordnet
sie der tatsächlichen 16-mal-4-Bank zu. Siehe
\href{https://arxiv.org/abs/hep-th/0109116}{Nath und Syed, vollständige
Spinorkopplungen}.

Aus dieser Zerlegung folgt weder, dass alle Bilineare verfügbare
Kontrollen sind, noch, dass jede Darstellung ein propagierendes Feld
ist. Produkte einfacher Matrizen im Einteilchenraum dürfen insbesondere
nicht automatisch als dieselben Produkte ihrer
Vielteilchen-Hamiltonoperatoren gelesen werden.

\subsection{6. Welche Operationen wirklich verfügbar sein
müssten}\label{welche-operationen-wirklich-verfuxfcgbar-sein-muxfcssten}

\subsubsection{6.1 Die algebraische Frage ist präziser
beantwortet}\label{die-algebraische-frage-ist-pruxe4ziser-beantwortet}

Im N=3-Sektor bleiben für verschiedene \textbf{angenommene}
Operationsalphabete folgende komplexe Kommutantdimensionen:

{\def\LTcaptype{none} % do not increment counter
\begin{longtable}[]{@{}lr@{}}
\toprule\noalign{}
Angenommenes Alphabet & Kommutantdimension \\
\midrule\noalign{}
\endhead
\bottomrule\noalign{}
\endlastfoot
X und Nb & 1444233216 \\
Zusätzlich dokumentierter Clock & 240742144 \\
Stattdessen zusätzlich SU(4)-Generatoren & 2247168 \\
Stattdessen zusätzlich Spin(10)-Generatoren & 1648 \\
Zusätzlich beide vollständigen Gruppen & 7 \\
Volle Gruppen und alle Modenbesetzungen & 1 \\
\end{longtable}
}

Die physische Verfügbarkeit ist somit nicht auf eine binäre Alternative
zusammengeschrumpft. Es gibt geprüfte Zwischenstufen, darunter auch die
in v1.6.6 behandelten einzelnen Casimir-Auslesungen.

Die volle Algebra nach Hinzunahme aller Besetzungen wurde unabhängig mit
einem kleineren Beweis kontrolliert. Der Einteilchen-Supportgraph der
gewährten Gruppengeneratoren ist verbunden, ebenso der Bosongraph. Ein
verbundener Graph hat verbundene Besetzungsgraphen mit k
ununterscheidbaren Fermionen, solange 0 kleiner k kleiner 64. Der native
Tensor verbindet die Fermion- und Bosonseiten. Die Besetzungsoperatoren
trennen alle Basiszustände. Polynomiale Spektralprojektoren und die
nichtverschwindenden Verbindungseinträge erzeugen daher alle
Matrixeinheiten.

So folgt bedingt die volle Algebra in N=2 und N=3. Gewährt man
zusätzlich geladene f-Instrumente, verbinden sich auch die beiden
Sektoren. Das ist eine vollständige \textbf{assoziativ-algebraische}
Aussage, kein Nachweis beliebiger unitärer Steuerbarkeit, effizienter
Messung oder nativer Instrumentenherkunft.

\subsubsection{6.2 Warum bloßes Lesen noch keine fehlende Bewegung
erzeugt}\label{warum-blouxdfes-lesen-noch-keine-fehlende-bewegung-erzeugt}

Ohne die zusätzlich gewährten Gruppenkontrollen erhalten X, Nb, der
Clock und alle Modenbesetzungen weiterhin mindestens eine nichtskalare
Farb-Cartanladung. Jede Komposition dieser Operatoren behält diese
Erhaltung. Die alleinige Erlaubnis, mehr Moden zu unterscheiden, hebt
das Hindernis also nicht auf.

Das ist die präzise Grenze zwischen \textbf{Auslesen} und
\textbf{Eingreifen}. Ein mathematisch definierter Projektor ist nicht
schon ein physisch ausführbares Messinstrument. Dass das Modell unter
einer Transformation symmetrisch ist, bedeutet nicht, dass es diese
Transformation als kontrollierten Eingriff erzeugt. Die bedingten
Zwei-Banken- und Vierzustands-Transferbeispiele aus v1.6.5 und v1.6.6
bleiben nützlich, aber ihre zusätzlichen Links bleiben auch zusätzliche
Voraussetzungen.

\subsection{7. Feldwörterbuch: gesicherte Teile und korrigierte
Reichweite}\label{feldwuxf6rterbuch-gesicherte-teile-und-korrigierte-reichweite}

Für die unveränderte antisymmetrische Paarmatrix und gleichhändige
Weylfelder verschwindet die ableitungsfreie skalare Kontraktion. Die
drei symmetrischen Spinortensoren tragen einen nichtverschwindenden
Kanal des Typs (1,0), mit konjugiertem Typ (0,1). Das ist am gesamten
ursprünglichen Tensor neu geprüft.

In der Konvention des gelieferten Programms hat jeder der drei Kanäle
pro Zeile Normquadrat 64. Die Summe ist 192 und die direkte Summennorm
folglich \(8\sqrt3\), \textbf{nicht 24}. Die Summe von drei einzelnen
Normen darf nicht mit der Norm ihrer orthogonalen direkten Summe
verwechselt werden. Eine anders normierte Vertexdefinition ändert diese
Zahlen, nicht den Nullkanal.

Der Satz gilt unter den genannten Feld- und Ableitungsannahmen.
Zusätzliche unabhängige Komponenten oder Ableitungskopplungen sind
andere Modelle und nicht generell ausgeschlossen. Keine dieser
Möglichkeiten ist dadurch bereits aus der ursprünglichen Bank
hergeleitet.

Auch die Zerlegung eines Vermittler-mal-Weyl-Komposits in die
Lorentztypen (3/2,0) und (1/2,0) ist eine Darstellungsaussage. Der
Ranganteil 2/6 gleich 1/3 ist das Gewicht im maximal gemischten Zustand
dieses Sechserraums. Für allgemeine Zustände kann derselbe Projektor
Gewichte von null bis eins haben. Er liefert nicht automatisch ein
Drittel des nativen Polgewichts. Die frühere Ladungsunterscheidung
bleibt bestehen: Die EOM-Größe D trägt Ladung minus eins, das ältere
erzeugende Komposit chi-dagger plus drei.

\subsection{8. Die fundamentale Frage: was fehlt
wirklich?}\label{die-fundamentale-frage-was-fehlt-wirklich}

\subsubsection{8.1 Ein überprüfbarer
Unterbestimmtheitssatz}\label{ein-uxfcberpruxfcfbarer-unterbestimmtheitssatz}

Betrachte bei festem W dieselbe Familie \(H_{\Delta,g}=\Delta N_b+gX\).
Alle diese Hamiltonoperatoren besitzen die gleiche innere
Quellsymmetrie, dieselbe Ladungserhaltung und denselben kommutierenden
endlichen Clock. Dennoch unterscheiden sie sich dynamisch. Schon für
dieselbe gefüllte Referenz F ist die dimensionslose Momentenkombination

\[\adjustbox{max width=\linewidth}{$\displaystyle 
\frac{\mu_4\mu_2}{\mu_3^2}
=1+1396\left(\frac{g}{\Delta}\right)^2.
$}\]

Bei g geteilt durch Delta gleich 1/20 ist sie 4,49, bei 1/40 dagegen
1,8725. Beide Parameterpunkte liegen im früher abgesicherten
Grundzustandsintervall. Weil die Größe dimensionslos ist, handelt es
sich nicht lediglich um eine Änderung der Maßeinheit. Die gleiche
Grammatik und Symmetrie lässt also unterschiedliche Antworten auf
dieselbe Präparation zu.

\textbf{Damit ist die eindeutige physische Dynamik aus diesen Grunddaten
allein nicht ableitbar.} Das widerlegt nicht die Möglichkeit einer
einfachen zusätzlichen Ursprungsregel. Es identifiziert genau deren
Aufgabe: Sie muss den Prozess samt verfügbarem Operationssatz und
Zustand auswählen, statt nur seine Symmetrie zu benennen. Im fest
gesetzten Prüfvertrag sind g und Delta natürlich definiert; offen ist
ihre physische Herkunft.

\subsubsection{8.2 Der kleinste tragfähige
Forschungsansatz}\label{der-kleinste-tragfuxe4hige-forschungsansatz}

Für eine einzelne Antwortrechnung sollte der relevante Raum nicht von
vornherein der gesamte Fockraum sein. Der zyklische Raum der
tatsächlichen Operationen auf dem angegebenen Startzustand ist der
kleinere Arbeitsraum. Die H-Lanczos-Kette ist eine exakte Darstellung
eines solchen Raums. Mehrere Präparationen, Messungen oder
Compileroperationen können jedoch einen größeren gemeinsamen Prozessraum
verlangen. Ihre Kompatibilität muss dann nachgewiesen werden; getrennte
kleine Modelle dürfen nicht einfach zusammengefügt werden. Eine
vollständige Rekonstruktion des Grundzustands aus dieser Referenz
verlangt außerdem einen Nachweis ihrer nichtverschwindenden Überlappung
mit Omega; die Ritz-Obergrenze allein ersetzt ihn nicht.

Die nächste physikalische Schließung bleibt dieselbe konkrete Bedingung:
Auf einer gemeinsam hergeleiteten Quelle und demselben Zustand muss eine
verfügbare Operation eine geladene Anregung zwischen operational
bestimmten Teilen übertragen. Der Test muss die Mehrzeitantwort von
bloßem Überlapp unterscheiden und ein konsistentes relativistisches
Feldwörterbuch verwenden. Koordinatenwechsel, Gram-Überlapp oder ein
innerer Clock allein erfüllen diese Bedingung nicht.

\subsection{9. Aktueller Stand T1 bis
T8}\label{aktueller-stand-t1-bis-t8}

Die folgende Tabelle ordnet die Beiträge dieser Runde den bisherigen
offenen Beweispflichten zu; sie benennt keine neue Definition der
ursprünglichen Tore.

{\def\LTcaptype{none} % do not increment counter
\begin{longtable}[]{@{}
  >{\raggedright\arraybackslash}p{(\linewidth - 4\tabcolsep) * \real{0.3333}}
  >{\raggedright\arraybackslash}p{(\linewidth - 4\tabcolsep) * \real{0.3333}}
  >{\raggedright\arraybackslash}p{(\linewidth - 4\tabcolsep) * \real{0.3333}}@{}}
\toprule\noalign{}
\begin{minipage}[b]{\linewidth}\raggedright
Tor
\end{minipage} & \begin{minipage}[b]{\linewidth}\raggedright
Beitrag dieser Runde
\end{minipage} & \begin{minipage}[b]{\linewidth}\raggedright
Weiter fehlender entscheidender Nachweis
\end{minipage} \\
\midrule\noalign{}
\endhead
\bottomrule\noalign{}
\endlastfoot
T1 & Exakte bedingte Kontrollalgebra; Auslese- und Eingriffslücke
präzisiert & Tatsächlich erzeugbare Compileroperationen und
Instrumente \\
T2 & Bessere lokale Struktur und Antwortgrenze in der endlichen Bank &
Quellseitiges Half-Charge-Feld mit Energie- und Adjungiertenkontrolle
sowie E8-/Clock-Zuordnung \\
T3 & Kleiner gemeinsamer Rechenraum innerhalb einer Bank & Gemeinsamer
physischer 3+1D-Ursprung \\
T4 & Vollständige innere Bilineare; korrekte Feldkanäle & Vollständiges
chirales Maß mit dynamischen Eichfreiheitsgraden \\
T5 & Exakter lokaler Hamilton-Lanczos-Anfang & Kontrollierter
relativistischer Kontinuums- und Streuungsgrenzfall \\
T6 & Expliziter Nachweis, dass die bisherigen Grunddaten g/Delta nicht
auswählen & Physische Kopplungs-, Skalen- und Spektralzuordnung,
einschließlich offener Neutrinofragen \\
T7 & Keine neue Schließung & Aus derselben Quelle erzeugter dynamischer
Spin-2-Sektor und universelle Kopplung \\
T8 & Unterscheidung Startzustand, Grundzustand und Ausleseinstrument
verschärft & Gemeinsame physische Präparation, Aufzeichnung und
Zustandsauswahl \\
\end{longtable}
}

Alle acht Tore bleiben offen. Die heutigen Ergebnisse schließen
begrenzte mathematische Fragen innerhalb des Programms, nicht diese
physikalischen Gesamtnachweise. Frühere RH-, Faktorisierungs- oder
P-versus-NP-Grenzen werden durch die endlichen Rechnungen ebenfalls
nicht verändert. Der begrenzte RH-Registerabgleich und seine
Aktualitätsgrenze sind in Teil B dokumentiert; ein vollständiger
Neuaudit dieser Arbeitsfronten fand nicht statt.

